\documentclass{eptcs}
\usepackage{hyperref}
\usepackage{amsmath, amssymb}
\usepackage{stmaryrd}
\usepackage{tikzit}
\newcommand{\tikzroot}{../tikz}
\IfFileExists{\tikzroot/zxw_styles.tikzdefs}{}{\renewcommand{\tikzroot}{extended-abstract/../tikz}}
\newcommand{\refsroot}{../refs}
\IfFileExists{\refsroot.bib}{}{\renewcommand{\refsroot}{extended-abstract/../refs}}
\input{\tikzroot/zxw_styles.tikzdefs}
\input{\tikzroot/zxw_styles.tikzstyles}
\usepackage[english]{babel}
\usepackage{enumitem}
\usepackage{float}

\newcommand{\Control}{\textsf{Ctrl}}
\newcommand{\Rule}[1]{\textsc{#1}}
\newcommand{\CSzero}{\Rule{CS0}}
\newcommand{\CScpy}{\Rule{CScpy}}
\newcommand{\CMzero}{\Rule{CM0}}
\newcommand{\CMcpy}{\Rule{CMcpy}}
\newcommand{\CMcom}{\Rule{CMcom}}
\newcommand{\csring}{\tilde{\mathcal{S}}_n}
\newcommand{\polyring}{\mathbb{C}[x_1,\ldots,x_n]/(x_1^2,\ldots,x_n^2)}
\newcommand{\lowerbox}[2][5]{\raisebox{-#1pt}{#2}}\newcommand{\wspids}[1][0.4]{\begin{tikzpicture}[scale=#1]
	\begin{pgfonlayer}{nodelayer}
		\node [style=uw] (0) at (0, 0) {};
		\node [style=none] (1) at (0.5, -0.5) {};
		\node [style=none] (2) at (-0.5, -0.5) {};
		\node [style=none] (3) at (0, 0.75) {};
    \node [style=label] (4) at (0, -0.5) {$\ldots$};
	\end{pgfonlayer}
	\begin{pgfonlayer}{edgelayer}
		\draw (0) to (1.center);
		\draw (2.center) to (0);
		\draw (3.center) to (0);
	\end{pgfonlayer}
\end{tikzpicture}}
\newcommand{\zspids}[1][0.4]{\begin{tikzpicture}[scale=#1]
  \begin{pgfonlayer}{nodelayer}
    \node [style=gn] (0) at (0, 0) {};
    \node [style=none] (1) at (0.5, -0.5) {};
    \node [style=none] (2) at (-0.5, -0.5) {};
    \node [style=none] (3) at (0, 0.75) {};
    \node [style=label] (4) at (0, -0.5) {$\ldots$};
  \end{pgfonlayer}
  \begin{pgfonlayer}{edgelayer}
    \draw (0) to (1.center);
    \draw (2.center) to (0);
    \draw (3.center) to (0);
  \end{pgfonlayer}
\end{tikzpicture}}
\newcommand{\numberstate}[1][a]{\begin{tikzpicture}[scale=0.4]
  \begin{pgfonlayer}{nodelayer}
    \node [style=gbox] (0) at (0, 0) {$#1$};
    \node [style=none] (1) at (0, 1) {};
  \end{pgfonlayer}
  \begin{pgfonlayer}{edgelayer}
    \draw (1.center) to (0);
  \end{pgfonlayer}
\end{tikzpicture}}
\newcommand{\numbergate}[1][a]{\begin{tikzpicture}[scale=0.5]
	\begin{pgfonlayer}{nodelayer}
		\node [style=gbox] (0) at (0, 0) {$#1$};
		\node [style=none] (3) at (0, 0.75) {};
		\node [style=none] (4) at (0, -0.75) {};
	\end{pgfonlayer}
	\begin{pgfonlayer}{edgelayer}
		\draw (3.center) to (0);
		\draw (0) to (4.center);
	\end{pgfonlayer}
\end{tikzpicture}}
\newcommand{\idwire}[1][0.5]{\begin{tikzpicture}[scale=#1]
	\begin{pgfonlayer}{nodelayer}
		\node [style=none] (0) at (0, 0) {};
		\node [style=none] (1) at (0, -1) {};
	\end{pgfonlayer}
	\begin{pgfonlayer}{edgelayer}
		\draw (0.center) to (1.center);
	\end{pgfonlayer}
\end{tikzpicture}}
\newcommand{\swap}[1][0.25]{\begin{tikzpicture}[scale=#1]
	\begin{pgfonlayer}{nodelayer}
		\node [style=none] (0) at (-1, -1) {};
		\node [style=none] (1) at (1, 1) {};
		\node [style=none] (2) at (1, -1) {};
		\node [style=none] (3) at (-1, 1) {};
	\end{pgfonlayer}
	\begin{pgfonlayer}{edgelayer}
		\draw [in=90, out=-90, looseness=0.75] (3.center) to (2.center);
		\draw [in=-90, out=90, looseness=0.75] (0.center) to (1.center);
	\end{pgfonlayer}
\end{tikzpicture}}
\newcommand{\coWs}[1][0.4]{\begin{tikzpicture}[scale=#1]
	\begin{pgfonlayer}{nodelayer}
		\node [style=dw] (0) at (0, 0) {};
		\node [style=none] (1) at (0.5, 0.5) {};
		\node [style=none] (2) at (-0.5, 0.5) {};
		\node [style=none] (3) at (0, -0.75) {};
		\node [style=label] (4) at (0, 0.6) {$\ldots$};
	\end{pgfonlayer}
	\begin{pgfonlayer}{edgelayer}
		\draw (0) to (1.center);
		\draw (2.center) to (0);
		\draw (3.center) to (0);
	\end{pgfonlayer}
\end{tikzpicture}}
\newcommand{\redzeroup}{\begin{tikzpicture}
	\begin{pgfonlayer}{nodelayer}
		\node [style={rn}] (0) at (0, 0) {};
		\node [style=none] (1) at (0, 0.5) {};
	\end{pgfonlayer}
	\begin{pgfonlayer}{edgelayer}
		\draw (0) to (1.center);
	\end{pgfonlayer}
\end{tikzpicture}}
\newcommand{\greenzeroup}{\begin{tikzpicture}
	\begin{pgfonlayer}{nodelayer}
		\node [style={gn}] (0) at (0, 0) {};
		\node [style=none] (1) at (0, 0.5) {};
	\end{pgfonlayer}
	\begin{pgfonlayer}{edgelayer}
		\draw (0) to (1.center);
	\end{pgfonlayer}
\end{tikzpicture}}
\setlength{\textfloatsep}{8pt plus 1pt minus 2pt}
\setlength{\floatsep}{7pt plus 1pt minus 1pt}
\setlength{\intextsep}{7pt plus 1pt minus 1pt}
\setlength{\abovedisplayskip}{4pt}
\setlength{\belowdisplayskip}{4pt}
\setlength{\abovedisplayshortskip}{2pt}
\setlength{\belowdisplayshortskip}{2pt}
\setlist[itemize]{leftmargin=*,nosep}
\setcounter{secnumdepth}{0}

\title{Algebraic Structure of Quantum Controlled States and Operators}
\author{Edwin Agnew \and Lia Yeh \and Richie Yeung}
\def\titlerunning{PolyZXW Extended Abstract}
\def\authorrunning{Agnew, Yeh, Yeung}

\begin{document}
\maketitle

\begin{abstract}
We show that controlled diagrams in the ZXW-calculus admit rich algebraic structure. Controlled states form a commutative ring isomorphic to multilinear polynomials, while controlled square matrices form a non-commutative ring. These properties supply new rewrite rules that make factorisation of arbitrary qubit Hamiltonians possible inside a single graphical calculus. The higher-order map $\Control$ recovers familiar controlled gates and composes cleanly, enabling multi-control and black-box Hamiltonian reasoning without leaving the diagrammatic setting.
\end{abstract}

Quantum control is a universal primitive that underlies block encodings, LCUs and Hamiltonian simulation. Yet graphical calculi typically treat control as a gadget rather than a first-class algebraic object. We work in the qubit ZXW-calculus (complete for finite-dimensional linear maps) and add controlled states and controlled square matrices as generators. The goal is to characterise their algebraic laws and exploit them for equational reasoning.
These rewrite rules already support Hamiltonian-level reasoning in practice: they are used for Hamiltonian summation and exponentiation~\cite{shaikh2022sum}, and the copy rule (\CMcpy) found in this work has been used to derive a general form for two-body electronic Hamiltonians under arbitrary linear encodings~\cite{mcdowallrose2025f2q}.
In this calculus, completeness for matrix-valued polynomials gives a purely diagrammatic route to factorising arbitrary qubit Hamiltonians.
Unlike traditional quantum graphical calculi whose parameters are restricted to phases and generator arity, here we allow higher-order reasoning by black-boxing controlled states and controlled square matrices as new generators.
This enables powerful new graphical rules for reasoning about quantum control in a wide variety of settings, where the axiomatic rewrites hold on the basis of the new generators being controlled diagrams irrespective of their internal structure.
The same framework unifies ordinary quantum controlled gates, multiple-control constructions, and algebraic manipulation of controlled states in one calculus. Interpreting controlled states and operators as multilinear polynomials also exposes direct links to classical algebraic tooling for optimisation and structural analysis.

\paragraph{Controlled states and controlled square matrices (Def.~2.1, 2.2).}
In extension of qubit graphical calculi, we introduce two additional generators.
\vspace{-0.5em}
\par
\noindent
\begin{minipage}[H]{0.49\linewidth}
\begin{equation}
  \scalebox{0.8}{$\left\llbracket \quad \tikzfig{\tikzroot/con/c_sq_def} \;\;\; \right\rrbracket$ } ~=~ \bra{0} \otimes I + \bra{1} \otimes M ~=~ \begin{bmatrix} I & M \end{bmatrix}
  \label{eq:def-csq}
\end{equation}
\end{minipage}\hfill
\begin{minipage}[H]{0.49\linewidth}
\vspace{0pt}
\begin{equation}
 \scalebox{0.8}{$\left\llbracket \;\;\; \tikzfig{\tikzroot/con/c_state_def} \;\;\; \right\rrbracket$ } ~=~ \ket{\vec{0}}\bra{0} + \ket{\psi}\bra{1} = \left[\ket{\vec{0}} \quad \ket{\psi}\right]
\label{eq:def-cstate}
\end{equation}
\end{minipage}
These are called controlled diagrams because their interpretation necessarily implies that they satisfy the (\CSzero) and (\CMzero) rules in Fig.~\ref{fig:ctrl-ring}: Their action when the control qubit is in the state $\ket{0}$ is to output $\ket{0...0}$ for states, and the identity map for square matrices.

\paragraph{The quantum control higher-order map}
As a prelude to our main results, we define the higher-order map $\Control$ to establish the applicability of our calculus to the quantum circuit setting.
$\Control$ sends any square matrix to its controlled square matrix acting on one more qubit wire---the control qubit. It is defined as the composition of the controlled square matrix generator with an arity-3 Z spider.
For example, the CX gate can be recovered as $\Control(X)$:
\begin{equation}
    \scalebox{0.8}{\tikzfig{\tikzroot/con/Fbox_cx}}
\end{equation}
Our results verify that $\Control$ preserves sequential and parallel composition, and that nesting $\Control$ recovers standard multi-control gates (e.g., Toffoli from $\Control(\Control(X))$.
Therefore, standard notions of quantum control are captured by the more general (i.e. not imposing unitarity or Hermiticity) semantics of our controlled generators.

\paragraph{Arithmetic diagrams and a normal form for them.}
Per Def.~5.1, a ZXW diagram with a single input on top is \textbf{arithmetic} if it contains only  $\lowerbox{\idwire}$, $\lowerbox{\swap}$ wires; $\lowerbox[10]{\wspids}$, $\lowerbox[10]{\zspids}$, $\lowerbox{\coWs}$ nodes; and $\lowerbox{\numberstate}$ boxes.
We show that all arithmetic diagrams must be controlled states, and conversely all controlled states are realisable as arithmetic diagrams.
In particular, we present an efficient algorithm to rewrite any arithmetic diagrams, and therefore any controlled states, to \emph{polynormal form (PNF)}, where control structure is arranged canonically so coefficients of \(p(x_1,\ldots,x_n)\) appear in the standard lexicographical ordering (Fig.~\ref{fig:pnf}).
\begin{figure}[H]
  \centering
  \scalebox{0.9}{\tikzfig{\tikzroot/poly/pnf}}
  \vspace{0.2em}
  {\small
  $
  ~=~
  \begin{bmatrix}
    1 & a_0 \\
    0 & a_1 \\
    \vdots & \vdots \\
    0 & a_{2^n-1}
  \end{bmatrix}$}
  \caption{Polynormal form for controlled states in the arithmetic fragment of the ZW-calculus.}
  \label{fig:pnf}
\end{figure}
This normal form is similar to the original ZW completeness result from which the original ZX-calculus completeness result was translated~\cite{hadzihasanovic2017thesis,Hadzihasanovic2018zwzxcomplete}, but for the first time this is for \emph{controlled} states instead of simply states.
This rule set was given for Cartesian Distributive Categories in~\cite{wilson2021revderbool,wilson2023diffpolycirc}, which proved completeness for Boolean circuits. This therefore addresses the question they raised as to whether the same rules are complete for continuous-valued, rather than discrete-valued, rings.

\paragraph{Isomorphism to multilinear polynomials.}
Let \(\tilde{\mathcal{S}}_n\) be the set of all controlled \(n\)-partite states.
Given controlled states $\tilde{\psi}$ and $\tilde{\phi}$, we define their addition $\tilde{\psi} \boxplus \tilde{\phi}$ and multiplication $\tilde{\psi} \boxtimes \tilde{\phi}$:
\begin{equation}
    \scalebox{0.8}{\tikzfig{\tikzroot/con/cs_add_def}} \qquad\qquad\qquad\qquad \scalebox{0.8}{\tikzfig{\tikzroot/con/cs_times_def}}
\end{equation}
We show that \((\tilde{\mathcal{S}}_n,\boxplus,\boxtimes)\) defines a commutative ring: \(\boxplus\) defines an Abelian group, \(\boxtimes\) defines a commutative monoid, and \(\boxtimes\) distributes over \(\boxplus\).
Furthermore, we investigate the precise nature of the correspondence to the ring of multilinear polynomials \(\mathbb{C}[x_1,\ldots,x_n]/(x_1^2,\ldots,x_n^2)\). For states, this is a folklore result in algebraic complexity, thereby essential in the study of multipartite entanglement. We extend this study to \emph{controlled} states.

First, we define the map $\phi_n: \polyring \to \csring$, which maps an arbitrary multilinear polynomial $p(x_1, ..., x_n) = a_0 + a_1x_n + ... + a_{2^n-1}x_1x_2...x_n$ to the polynormal form in Fig.~\ref{fig:pnf}.
What we find is that the ring of controlled states is not only in bijection, but in \emph{isomorphism} to the ring of multilinear polynomials, where $\phi_n$ defines the isomorphism. In particular, the addition and multiplication operations on controlled states---identified with the W and Z generators respectively---are in one-to-one correspondence with the addition and multiplication of polynomials, and the rewrite rules for controlled states correspond to the algebraic properties of polynomials. This means that we can use the rich algebraic structure of polynomials to reason about controlled states in a powerful way, and that any algebraic manipulation of polynomials can be translated into a corresponding manipulation of controlled states in the ZXW-calculus.

\paragraph{Completeness for matrix-valued polynomials}
Having proven that controlled states form a ring isomorphic to multilinear polynomials, the question is what algebraic structure controlled square matrices have.
Let \(\tilde{M}_n\) be the set of all controlled square matrices on \(n\) qubits.
We define their weighted addition \(\widetilde{\sum_i c_iM_i}\) and multiplication \(\widetilde{\prod_i M_i}\):
\begin{equation}
  \scalebox{0.74}{\tikzfig{\tikzroot/con/csq_sum}}
\qquad\qquad\qquad
\scalebox{0.74}{\tikzfig{\tikzroot/con/csq_prod}}
\end{equation}
We give an efficient algorithm to write any controlled square matrix to polynormal form: Higher-order generators are not repeated, and control structure is arranged canonically so coefficients of \(p(x_1,\ldots,x_n)\) appear in PNF for controlled states (Fig.~\ref{fig:pnf}).
Moreover, we prove that \((\tilde{M}_n,+,\times)\) forms a non-commutative ring---isomorphic to multivariate polynomials over same-size square matrices with coefficients in $\mathbb{C}$---due to satisfying all ring axioms:
\begin{enumerate}[noitemsep]
  \item Scalar multiplication.\\ For \(c\in\mathbb{C}\), set \(\tilde{a}\times\tilde{M}:=\widetilde{aM}\), i.e. the multiplication defined in Eq.~(5) by $\lowerbox{\numberstate}$.
  \item Commutativity of addition.\\ This is axiomatised in (\CMcom).
  \item Additive identity.\\ The additive identity is defined as $\lowerbox[3]{\scalebox{0.8}{\redzeroup}} \otimes I_n$, as the unit of $\lowerbox[10]{\wspids}$ satisfying \(\tilde M + \tilde{0}=\tilde M\).
  \item Additive inverse.\\ The additive inverse of $\tilde M$ is \(\widetilde{-M}:= \raisebox{-5pt}{\numbergate[-1]} \circ \tilde{M}\), satisfying \(\tilde M + \widetilde{-M}=\tilde{0}\).
  \item Associativity of addition.\\ This follows from associativity of $\lowerbox[10]{\wspids}$: \((\tilde M_1 + \tilde M_2) + \tilde M_3=\tilde M_1 + (\tilde M_2 + \tilde M_3)\).
  \item Multiplicative identity.\\ Like the additive identity, the multiplicative identity is $\lowerbox[3]{\scalebox{0.8}{\greenzeroup}} \otimes I_n$, as the unit of $\lowerbox[10]{\zspids}$ satisfying \(\tilde M \times \tilde{1}=\tilde M\).
  \item Associativity of multiplication.\\ This follows from associativity of $\lowerbox[10]{\zspids}$: \((\tilde M_1\times\tilde M_2)\times\tilde M_3=\tilde M_1\times(\tilde M_2\times\tilde M_3)\).
  % \item Multiplicative inverse (when defined; not a ring axiom).\\ If \(M\) is invertible, then \(\tilde M\times\widetilde{M^{-1}}=\widetilde{I}=\widetilde{M^{-1}}\times\tilde M\).
  \item Distributivity of multiplication over addition.\\
  \item We prove this from (\CMcom), (\CMcpy), and (\Rule{BZW}). Thus, \(\tilde M_1 \times (\tilde M_2 + \tilde M_3)=(\tilde M_1 \times \tilde M_2) + (\tilde M_1 \times \tilde M_3)\):
  % We prove this follows from (\CMcom), (\CMcpy), and (\Rule{BZW}). Therefore, we have \(\tilde M_1 \times (\tilde M_2 + \tilde M_3)=(\tilde M_1 \times \tilde M_2) + (\tilde M_1 \times \tilde M_3)\):
  \begin{equation}
    \scalebox{0.8}{\tikzfig{\tikzroot/con/csq_dist_statement}}
  \end{equation}
\endgroup

\textbf{Summary of completeness results.}
\begin{figure}[H]
  \centering
  \renewcommand{\arraystretch}{1.25}
  \begin{tabular}{|p{0.45\textwidth}|p{0.45\textwidth}|}
    \hline
    \multicolumn{2}{|c|}{\textbf{Controlled Ring Rules (Subset of Fig.~1)}} \\
    \hline
    \begin{minipage}{\linewidth}
      \vspace{-1.2em}
      \begin{gather*}
        \tikzfig{\tikzroot/con/c_state0}\\[-0.4em]
        (\CSzero)\\
        \scalebox{0.9}{\tikzfig{\tikzroot/con/cs_copy}}\\[-0.4em]
        (\CScpy)
      \end{gather*}
    \end{minipage} &
    \begin{minipage}{\linewidth}
      \vspace{-0.9em}
      \begin{gather*}
        \scalebox{0.9}{\tikzfig{\tikzroot/con/c_sq0}}\\[-0.4em]
        (\CMzero)\\
        \scalebox{0.85}{\tikzfig{\tikzroot/con/csq_copy}}\\[-0.4em]
        (\CMcpy)\\
        \scalebox{0.85}{\tikzfig{\tikzroot/con/csq_add_comm_statement}}\\[-0.4em]
        (\CMcom)
      \end{gather*}
      \vspace{-1.7em}
    \end{minipage}\\
    \hline
  \end{tabular}
  \caption{In this work, we prove that adding these Controlled Ring Rules to the complete ZW- or ZXW-calculus rules attains completeness for all arithmetic diagrams (Theorem~5.1), completeness for all ring axioms for controlled states (Theorem~5.2), and completeness for all ring axioms for controlled square matrices (Theorem~6.1).}
  \label{fig:ctrl-ring}
\end{figure}

\begingroup
\footnotesize
\bibliographystyle{abbrv}
\bibliography{\refsroot}
\endgroup

\end{document}
